\documentclass[../../main.tex]{subfiles}
\begin{document}

{\bf [解]}

\paragraph{(1)}

% TODO: n=1,n=2の時の遷移図
\begin{figure}[htb]
  \centering
  \begin{tikzpicture}[
      x=3.2cm, y=1.8cm, >=stealth,
      box/.style={draw=gray!70!black, fill=white, rounded corners=3pt, line width=0.8pt, inner sep=4pt, font=\small},
      state/.style={font=\footnotesize\bfseries},
      prob/.style={font=\scriptsize, fill=white, inner sep=1.5pt}
    ]

    % ステップ見出し
    \node[font=\footnotesize\bfseries, gray!60!black] at (0, 0.9) {初期状態};
    \node[font=\footnotesize\bfseries, gray!60!black] at (1, 0.9) {1回目終了時};
    \node[font=\footnotesize\bfseries, gray!60!black] at (2, 0.9) {2回目};

    % 初期ノード (白 a+2, 赤 1, 合計 a+3)
    \node[box] (init) at (0, 0) {
      \begin{tabular}{cc}
        白 & 赤 \\ \hline
        $a+2$ & $1$
      \end{tabular}
    };

    % 1回目分岐後のノード
    % 上: 赤を引いた場合 (確率 1/(a+3) = p_1)
    \node[box, draw=red!80!black] (s1_red) at (1, 0.5) {
      \begin{tabular}{c}
        白 \\ \hline
        $a+2$
      \end{tabular}
    };
    % 下: 白を引いた場合 (確率 (a+2)/(a+3))
    \node[box, draw=blue!80!black] (s1_white) at (1, -0.5) {
      \begin{tabular}{cc}
        白 & 赤 \\ \hline
        $a$ & $1$
      \end{tabular}
    };

    % 1回目の遷移矢印
    \draw[->, thick, red!85!black] (init) -- (s1_red)
    node[midway, prob, above left] {赤：$p_1 = \frac{1}{a+3}$};
    \draw[->, thick, blue!85!black] (init) -- (s1_white)
    node[midway, prob, below left] {白：$\frac{a+2}{a+3}$};

    % 2回目の結果ノード
    % s1_red から: 赤玉なし (確率 0)
    \node[draw=gray!50, rounded corners=2pt, font=\footnotesize, text=gray!60!black] (s2_zero) at (2, 0.5) {赤は引けない (確率 $0$)};

    % s1_white から: 赤を引く (確率 1/(a+1))
    \node[draw=red!85!black, fill=red!8, rounded corners=3pt, line width=1.0pt, font=\footnotesize\bfseries] (s2_hit) at (2, -0.2) {赤玉を取り出す};
    \node[draw=gray!50, rounded corners=2pt, font=\footnotesize, text=gray!60!black] (s2_miss) at (2, -0.8) {白玉を取り出す};

    % 2回目の遷移矢印
    \draw[->, dashed, gray!60] (s1_red) -- (s2_zero)
    node[midway, prob, above] {確率 $0$};
    \draw[->, line width=1.1pt, red!85!black] (s1_white) -- (s2_hit)
    node[midway, prob, above left] {赤：$\frac{1}{a+1}$};
    \draw[->, gray!60] (s1_white) -- (s2_miss)
    node[midway, prob, below left] {白：$\frac{a}{a+1}$};

  \end{tikzpicture}
  \caption{$n=1, 2$ における袋の状態推移と赤玉を取り出す確率}
  \label{fig:transition_tree_n1_n2}
\end{figure}

$n=1$の時，最初は白玉が$a+2$個，赤玉が$1$個だから赤玉を取り出す確率は
\begin{align}
  p_1 = \dfrac{1}{a+3} \label{eq:1}
\end{align}
である．

$n=2$の時，$1$回目に白玉を引けば袋の中身は白が$a$個，赤が$1$個であり，
$1$回目に赤玉を引けば袋の中身は白が$a+2$個，赤玉が$0$個だから，
赤玉を取り出す確率は
\begin{align}
  p_2 = (1-p_1)\dfrac{1}{a+1} + p_1\cdot 0 =  \dfrac{a+2}{(a+3)(a+1)} \label{eq:2}
\end{align}
である．
\cref{eq:1,eq:2}をまとめて
\begin{align}
  p_1 = \dfrac{1}{a+3}, p_2 = \dfrac{a+2}{(a+3)(a+1)} \label{eq:3}
\end{align}
である．

\paragraph{(2)}

% TODO: 一般の場合の遷移の様子
\begin{figure}[htb]
  \centering
  \begin{tikzpicture}[
      x=3.8cm, y=1.8cm, >=stealth,
      box/.style={draw=blue!85!black, fill=white, rounded corners=4pt, line width=1.0pt, inner sep=6pt, font=\small},
      prob/.style={font=\footnotesize\bfseries, fill=white, inner sep=2pt}
    ]

    % 列ヘッダー
    \node[font=\normalsize\bfseries, gray!70!black] at (0, 1.1) {$n$ 回目終了時の袋の状態};
    \node[font=\normalsize\bfseries, gray!70!black] at (1.5, 1.1) {$n+1$ 回目の試行};

    % 状態 [ア]: 白 a, 赤 1 (n回目に白を引いた: 確率 1 - p_n)
    \node[box] (stateA) at (0, 0.2) {
      \begin{tabular}{c}
        \textbf{状態 [ア]} \ （直前が白）\\[2pt]
        \begin{tabular}{|cc|}
          \hline
          白 & 赤 \\ \hline
          $a$ & $1$ \\ \hline
        \end{tabular}\\[4pt]
        確率：$1 - p_n$
      \end{tabular}
    };

    % 状態 [イ]: 白 a, 赤 0 (n回目に赤を引いた: 確率 p_n)
    \node[box, draw=gray!60] (stateB) at (0, -1.4) {
      \begin{tabular}{c}
        \textbf{状態 [イ]} \ （直前が赤）\\[2pt]
        \begin{tabular}{|c|}
          \hline
          白 \\ \hline
          $a$ \\ \hline
        \end{tabular}\\[4pt]
        確率：$p_n$
      \end{tabular}
    };

    % 次のステップの目標: 赤を引く (p_{n+1})
    \node[draw=red!85!black, rounded corners=4pt, line width=1.2pt, inner sep=6pt, font=\normalsize\bfseries]
    (nextRed) at (1.5, 0.4) {赤玉を取り出す（$p_{n+1}$）};

    % 次のステップ: 白を引く
    \node[draw=gray!50, rounded corners=3pt, font=\small, text=gray!60!black]
    (nextWhite) at (1.5, -0.3) {白玉を取り出す};

    % 状態 [ア] からの遷移
    \draw[->, line width=1.4pt, red!85!black] (stateA) -- (nextRed)
    node[midway, above, red!90!black] {確率 $\frac{1}{a+1}$};
    \draw[->, line width=0.9pt, gray!60] (stateA) -- (nextWhite)
    node[midway, below left] {確率 $\frac{a}{a+1}$};

    % 状態 [イ] からの遷移 (赤玉がないので確率 0)
    \draw[->, line width=1.0pt, dashed, gray!50] (stateB.east) -- (nextRed.south west)
    node[midway, below right, gray!60!black] {確率 $0$ （赤玉なし）};

  \end{tikzpicture}
  \caption{一般の $n$ 回目から $n+1$ 回目への状態遷移と確率漸化式の関係}
  \label{fig:general_step_markov_transition}
\end{figure}

$n \ge 2$とする．$n$回目の操作後の箱の状態は以下のいずれかである．
\begin{itemize}
  \item[ア] 白$a$個，赤$1$個 （$n$回目に白を引く）
  \item[イ] 白$a$個 （$n$回目に赤を引く）
\end{itemize}

$n+1$回目に赤を引くのは，[ア]の場合の時に赤を引く場合だから
\begin{align}
  p_{n+1} = \dfrac{1}{a+1} (1-p_n) \quad (n \ge 2)
\end{align}
である．変形して
\begin{align}
  p_{n+1} - \dfrac{1}{a+2} = \dfrac{-1}{a+1} \left( p_n - \dfrac{1}{a+2} \right)
\end{align}
である．漸化式をくり返し用いて，$n \ge 2$のとき
\begin{align}
  p_n
  &= \left(-\dfrac{1}{a+1}\right)^{n-2} \left( p_2 - \dfrac{1}{a+2} \right) + \dfrac{1}{a+2}
\end{align}
であり，\cref{eq:3}を代入して
\begin{align}
  p_n
  &= (-1)^{n-2} \left( \dfrac{1}{a+1} \right)^{n-1} \dfrac{1}{(a+2)(a+3)} + \dfrac{1}{a+2} \\
  &= \dfrac{1}{a+2} \left\{ 1 - \left(-\dfrac{1}{a+1}\right)^{n-1} \dfrac{1}{a+3} \right\} \label{eq:4}
\end{align}
である．

\paragraph{(3)}

\cref{eq:4}は\cref{eq:3}より$n=1$でも成立する．
従って$A_m = \sum_{n=1}^m p_n$ とすると，
\begin{align}
  A_m
  &= \dfrac{1}{a+2}\sum_{n=1}^m \left\{ 1 - \dfrac{1}{a+3} \left(-\dfrac{1}{a+1}\right)^{n-1} \right\} \\
  &= \dfrac{1}{a+2}\left[m - \dfrac{1}{a+3} \dfrac{1 - \left(-\dfrac{1}{a+1}\right)^m}{1 - \left(-\dfrac{1}{a+1}\right)}\right] \\
  &= \dfrac{1}{a+2}\left[m - \dfrac{1}{a+3} \dfrac{a+1}{a+2} \left\{ 1 - \left(-\dfrac{1}{a+1}\right)^m \right\}\right]
\end{align}
である．両辺を$m$で割って，$a>0$に注意して求める極限値は
\begin{align}
  \dfrac{1}{m} A_m
  & = \dfrac{1}{a+2} - \dfrac{1}{m} \dfrac{a+1}{(a+2)(a+3)} \left\{ 1 - \left(-\dfrac{1}{a+1}\right)^m \right\} \xrightarrow{m\to\infty} \dfrac{1}{a+2}
\end{align}
である．

\end{document}
